% !TeX root = main.tex
\changecolours{ScoutNavy}
\section[CRUD Avançado: Filtros e Atualizações Atômicas]{CRUD Avançado: Filtros Complexos e Atualizações Atômicas}

% ------------------------------------------------------------------------------
\begin{frame}[fragile]{O Operador \texttt{\$elemMatch} e o Problema dos Falsos Positivos}
  \vspace{0.05cm}
  \begin{columns}[t]
    \begin{column}{0.48\textwidth}
      \alert{A Armadilha do Filtro Ingênuo em Arrays}\par
      \itemseps[topsepi=0.2mm,itemsepi=0.5mm,labelsep=1.5mm,leftmargini=2.5mm]
      \scriptsize
      Considere o pedido com dois itens:
      \begin{itemize}
        \item Item A: \texttt{preco = 250}, \texttt{qtd = 1}
        \item Item B: \texttt{preco = 20}, \texttt{qtd = 8}
      \end{itemize}
      Se consultarmos:
      \begin{lstlisting}[language=mongosh]
db.pedidos.find({
  "itens.preco": { $gt: 100 },
  "itens.qtd": { $gte: 5 }
})
      \end{lstlisting}
      \textbf{Resultado:} O documento é retornado! O Item A atendeu ao preço e o Item B à quantidade. O filtro avaliou itens \textbf{diferentes}!
    \end{column}
    \begin{column}{0.48\textwidth}
      \alert{A Solução Estrita com \texttt{\$elemMatch}}\par
      \scriptsize
      O operador \texttt{\$elemMatch} força o motor a exigir que \textbf{todas as condições sejam atendidas no mesmo elemento}:
      \begin{lstlisting}[language=mongosh]
db.pedidos.find({
  itens: {
    $elemMatch: {
      preco: { $gt: 100 },
      qtd: { $gte: 5 }
    }
  }
})
      \end{lstlisting}
      \textbf{Resultado:} Nenhum documento é retornado, eliminando com precisão o falso positivo!
    \end{column}
  \end{columns}
\end{frame}

% ------------------------------------------------------------------------------
\begin{frame}[fragile]{Operador \texttt{\$expr}: Comparação de Campos Internos}
  \vspace{0.05cm}
  \begin{columns}[t]
    \begin{column}{0.48\textwidth}
      \alert{Quando Usar o \texttt{\$expr}?}\par
      \itemseps[topsepi=0.2mm,itemsepi=0.5mm,labelsep=1.5mm,leftmargini=2.5mm]
      \scriptsize
      \begin{itemize}
        \item Filtros comuns comparam um campo com um valor estático (\texttt{preco > 100}).
        \item Para comparar \textbf{dois campos do mesmo documento}, usamos \texttt{\$expr}.
        \item O prefixo \texttt{\$} indica que o valor daquele campo deve ser resolvido em tempo de execução.
      \end{itemize}
    \end{column}
    \begin{column}{0.48\textwidth}
      \alert{Exemplos Práticos no E-Commerce}\par
      \begin{lstlisting}[language=mongosh]
// 1. Produtos com estoque atual abaixo do estoque minimo:
db.produtos.find({
  $expr: {
    $lt: ["$estoque.atual", "$estoque.minimo"]
  }
})

// 2. Pedidos em que total pago foi menor que o subtotal:
db.pedidos.find({
  $expr: {
    $lt: ["$pagamento.total_pago", "$pagamento.subtotal"]
  }
})
      \end{lstlisting}
    \end{column}
  \end{columns}
\end{frame}

% ------------------------------------------------------------------------------
\begin{frame}[fragile]{Atualizações Atômicas Concorrentes: \texttt{\$inc}, \texttt{\$min} e \texttt{\$max}}
  \vspace{0.05cm}
  \begin{columns}[t]
    \begin{column}{0.48\textwidth}
      \alert{Por que Evitar \texttt{\$set} para Contadores?}\par
      \itemseps[topsepi=0.2mm,itemsepi=0.5mm,labelsep=1.5mm,leftmargini=2.5mm]
      \scriptsize
      \begin{itemize}
        \item Ler o estoque para a memória da aplicação, subtrair $1$ e salvar com \texttt{\$set} causa \textbf{Condições de Corrida (\textit{Race Conditions})}.
        \item Se duas requisições concorrentes lerem estoque $10$, ambas gravarão $9$, perdendo uma venda!
        \item O operador \textbf{\texttt{\$inc}} opera diretamente no bloco de memória do WiredTiger com garantia de atomicidade.
      \end{itemize}
    \end{column}
    \begin{column}{0.48\textwidth}
      \alert{Sintaxe de Operadores Atômicos}\par
      \begin{lstlisting}[language=mongosh]
// Decrementar estoque de forma atomica:
db.produtos.updateOne(
  { sku: "SKU-INF-000101" },
  { $inc: { "estoque.atual": -2 } }
)

// Atualizar apenas se for o menor preco historico:
db.produtos.updateOne(
  { sku: "SKU-INF-000101" },
  { $min: { "historico.menor_preco": 299.90 } }
)
      \end{lstlisting}
    \end{column}
  \end{columns}
\end{frame}

% ------------------------------------------------------------------------------
\begin{frame}[fragile]{Manipulação de Arrays: \texttt{\$push}, \texttt{\$slice} e \texttt{arrayFilters}}
  \vspace{0.05cm}
  \begin{columns}[t]
    \begin{column}{0.48\textwidth}
      \alert{Manutenção do \textit{Subset Pattern}}\par
      \scriptsize
      Ao inserir avaliações ou logs, usamos \texttt{\$push} com o modificador \texttt{\$slice: -5} para manter o array limitado aos últimos 5 itens:
      \begin{lstlisting}[language=mongosh]
db.produtos.updateOne(
  { sku: "SKU-INF-000101" },
  {
    $push: {
      ultimas_avaliacoes: {
        $each: [{ nota: 5, autor: "Maria" }],
        $slice: -5 // Mantem apenas os 5 mais recentes
      }
    }
  }
)
      \end{lstlisting}
    \end{column}
    \begin{column}{0.48\textwidth}
      \alert{Modificação Seletiva com \texttt{arrayFilters}}\par
      \scriptsize
      Atualiza um subdocumento específico de um array usando o identificador \texttt{\$[alvo]}:
      \begin{lstlisting}[language=mongosh]
db.pedidos.updateOne(
  { codigo: "PED-2026-001" },
  {
    $set: {
      "itens.$[alvo].status_entrega": "SEPARADO",
      "itens.$[alvo].data": new Date()
    }
  },
  {
    arrayFilters: [ { "alvo.sku": "ITEM-A" } ]
  }
)
      \end{lstlisting}
    \end{column}
  \end{columns}
\end{frame}
